\section{Introduction}

Wall crossing for bundles on a K3 surface relates compact holomorphic
symplectic manifolds through contractions of exceptional subvarieties.
The local quotient describes those subvarieties and their intersections.
To calculate their attachments to the compact moduli space one must also
identify their integral classes in its global cohomology ring. The same
distinction arises for wall maps: a local change of exceptional
configuration does not specify the action of a compact correspondence on
universal classes.

We carry out these identifications for a rank-three family with three
central line factors and cross-extension dimensions $1,2,1$.
Its six nearby chambers have two types of deepest exceptional fibre:
two planes meeting transversely at a point, or a plane and the first
Hirzebruch surface meeting along its negative section. The compact family,
its primitive exceptional cycles, and its harmonic metric degeneration
were constructed in \cite{wuebben17}. That calculation determined a
middle-degree attachment with the marked complement ring as an input.
The general descent statement in \cite{wuebben18} likewise retains the
geometric overlaps. Here we calculate the complement ring and all its
ordinary integral cohomological attachment maps, and identify the
correspondences induced by the actual compact slope walls.

\subsection{The family and the main result}
Fix a marked projective K3 surface $X$ with N\'eron--Severi lattice
generated by $H,D_1,D_2$, where $H$ is ample and the intersection matrix is
\begin{equation}\label{eq:NS}
 \begin{pmatrix}2&0&0\\0&-6&-4\\0&-4&-8\end{pmatrix}.
\end{equation}
The marking of $H^2(X,\Z)=\Lambda=3U\oplus2E_8(-1)$ is chosen so that
\begin{equation}\label{eq:embedding}
 H=e_1+f_1,\qquad D_1=e_2-3f_2,\qquad
 D_2=-4f_2+e_3-4f_3,
\end{equation}
where $e_i\cdot f_i=1$ in the three hyperbolic planes. Put
\begin{equation}\label{eq:vectors}
 \begin{gathered}
 a_0=(1,0,1),\quad a_1=(1,D_1,-2),\quad a_2=(1,D_2,-3),\\
 v=a_0+a_1+a_2=(3,D_1+D_2,-4),\qquad w=a_0+a_2.
 \end{gathered}
\end{equation}
Vectors are written in the Mukai convention $v(F)=\ch(F)\sqrt{\td_X}$,
with pairing $(r,c,s)\cdot(r',c',s')=c\cdot c'-rs'-r's$.
Let $M_\sigma$ be the compact moduli space of slope-stable bundles of
vector $v$ in a sufficiently small generic chamber of
$H+sD_1+tD_2$. Its central classifying map is denoted
$\pi_\sigma:M_\sigma\to Z$. The geometric inputs defining this map are
restated in Section~\ref{sec:geometry}. Let $o\in Z$ be the class of
$\OO_X\oplus\OO_X(D_1)\oplus\OO_X(D_2)$, and set
\[
 E_\sigma=(\pi_\sigma^{-1}(o))_{\mathrm{red}},\qquad
 V_\sigma=M_\sigma\setminus E_\sigma.
\]
The three wall coordinates are
\begin{equation}\label{eq:levels}
 \zeta(s,t)=\frac{2\pi}{3}
 \begin{pmatrix}-10&-12\\8&0\\2&12\end{pmatrix}
 \binom{s}{t},\qquad \zeta_0+\zeta_1+\zeta_2=0.
\end{equation}
Chamber $A$ has signs $(+,-,+)$.

For the integral answer let $Y=X^{[2]}$ and $R=H^*(Y,\Z)$, with its
complex orientation. Write $\mu:\Lambda\to R^2$ for the usual
Hilbert-square map, $\delta$ for half the effective Hilbert--Chow
divisor, and $A\in R^4$ for the class of subschemes containing a fixed
point. Section~\ref{sec:ring} gives a finite integral basis and all
products in $R$. Define
\begin{equation}\label{eq:planes-intro}
 \begin{aligned}
 u_0&=7H-2D_1-2D_2,& u_2&=-5H+2D_1+D_2,\\
 p_0&=A+3\delta^2+\tfrac12\bigl(\mu(u_0)^2+5\delta\mu(u_0)\bigr),\\
 p_2&=A+\delta^2+\tfrac12\bigl(\mu(u_2)^2-3\delta\mu(u_2)\bigr).
 \end{aligned}
\end{equation}
The classes $p_0,p_2$ are integral, although these expressions use halves.
The universal families will be normalized by the rank-one index of
$y=(0,-f_2,0)$; the definition is given in
\eqref{eq:normalized-index}.

\begin{theorem}\label{thm:main}
For the marked compact family above the following statements hold.
\begin{enumerate}
\item The four adjacent crossings on $\zeta_0=0$ and $\zeta_2=0$ are
elementary Mukai flops. The two crossings on $\zeta_1=0$ are
isomorphisms over the corresponding partial wall target. In the Mukai
markings their integral ring actions are respectively the identity and
the reflection $s_{a_1}(z)=z+(z,a_1)a_1$. For a flop the ring action is
induced by $[\Gamma]+[P\times P']$, where $\Gamma$ is the graph closure
and $P,P'$ are the exceptional planes.
\item The explicit Mukai isometry
\[
 g=e^{2H-D_1-D_2}s_{(1,H,2)},\qquad g(v)=(1,0,-1),
\]
specifies oriented integral ring identifications
$\Phi_\sigma:H^*(M_\sigma,\Z)\to R$. In these identifications
\begin{equation}\label{eq:complement-main}
 H^*(V_\sigma,\Z)=R/(p_0,p_2).
\end{equation}
The groups are free of ranks $1,23,274,21$ in degrees $0,2,4,6$ and
zero otherwise. All determinant-normalized universal-index Chern classes
are given by the topological formula in
Proposition~\ref{prop:indices} and reduction modulo $(p_0,p_2)$.
\item Let $W_\sigma$ be a regular neighborhood of $E_\sigma$ and give
$L_\sigma=\partial W_\sigma$ its outward orientation. In chamber $A$,
the degree-four restriction to $H^4(L_A,\Z)=\Z/8$ is
\begin{equation}\label{eq:kappa-intro}
 [z]\longmapsto\int_Yz(p_0-3p_2)\pmod8.
\end{equation}
Sections~\ref{sec:attachments}--\ref{sec:actions} give the neighborhood
and boundary rings, all restriction, relative and connecting maps in
each chamber, and the wall actions on their integral generators.
For the fixed normalizer $y$, the divisorial action on indices includes
the inverse determinant line of the $y-a_1$ index. The six-chamber
composition is the identity on integral cohomology.
\end{enumerate}
\end{theorem}

The first assertion supplies the geometric maps used in the other two.
At an isolated wall the exhaustive factor list leaves two rigid factors
with cross-extension dimension three. Their whole compact fibre is the
plane in the incidence resolution. At the divisorial wall, the second
factor varies over an entire smooth K3 surface, including the curve
lying above $o$. Both resolutions are the intrinsic blow-up of the
reduced singular surface in the partial wall target. A determinant
section has a simple zero on the exceptional divisor; its opposite
signs in the two Mukai markings force the reflection.

To calculate the integral maps we use the diagonal blow-up of $X\times X$
and classical integral Hilbert-square cohomology. Markman's universal
correspondence for the displayed $g$ identifies this calculation with
the chamber spaces. The exceptional classes then have the coordinates
\eqref{eq:planes-intro}. Primitive restriction witnesses identify the
kernel of restriction to $V_\sigma$ as their ideal. Poincar\'e--Lefschetz
duality computes the boundary restrictions and fixes their connecting
signs. This order of construction uses global K3 geometry before
excision; local singularity types alone are not inputs sufficient to
recover multiplication.

\subsection{Relation to earlier results}
The finite exceptional configurations already occur in
Bielawski--Dancer \cite[Example~6.10]{bd}. Integral cohomology of the
corresponding unimodular hypertoric quotients is covered by
Hausel--Sturmfels \cite[Theorem~1.1]{hs}; their core intersection forms
are treated by Hausel--Swartz \cite[Theorem~2]{hsw}.
Compact configurations of planes and ruled surfaces under successive
Mukai flops also occur in the K3 and EPW construction of
Donten-Bury et al.\ \cite[\S3.1 and Remark~6.17]{epw}.
Thus neither the two local fibre types nor the existence of interacting
compact exceptional surfaces is the assertion of this paper.

The integral basis used below is due to Qin--Wang
\cite[Theorem~5.4 and Remark~5.6]{qw}, in the geometric form of
Boissi\`ere--Nieper-Wi\ss kirchen--Sarti \cite[Remark~5.7]{bns}.
Markman's integral transport and universal-class compatibility
\cite{markman08} provide the comparison with $Y$.
The corrected flop correspondence is classical
\cite{riess,llw}; identity and reflection actions at K3 walls have
their general Bridgeland counterparts in
Bayer--Macr\`i \cite[Lemma~10.1]{bm}.
Our contribution is the compact slope-wall identification together
with the full integral calculation in its specified marking, including
the complement, the six attachments and the universal normalization.
The quotient-ideal argument and the cohomological six-cycle relation
are consequences of these identifications.

Finally Section~\ref{sec:metrics} pairs the calculated classes with the
harmonic K\"ahler class. This gives the signed line periods and ruled
surface volumes and identifies the cycles in the metric theorems of
\cite{wuebben17,wuebben18}. The analytic estimates are imported with
their hypotheses. We make no claim about Chow-level path relations,
monodromy of a complex smoothing, or reconstruction of an underlying
manifold or metric tensor from the cohomology ring.
